\section{Register File}

\subsection{What the register file is}

The register file is the CPU’s small, fast architectural storage. RISC-V RV32 has 32 integer registers:

\begin{lstlisting}[language={}]
x0, x1, x2, ..., x31
\end{lstlisting}

Each is 32 bits wide in RV32. Most instructions read one or two registers and write one destination register.

\subsection{The special rule for x0}

RISC-V register \texttt{x0} is hardwired to zero. Reading \texttt{x0} always returns zero. Writing \texttt{x0} has no effect. This simplifies many instructions. For example:

\begin{lstlisting}[language={}]
addi x5, x0, 123   # load small constant 123 into x5
add  x6, x7, x0    # copy x7 into x6
\end{lstlisting}

\texttt{rtl/regfile.v} implements this rule by returning zero when a read address is zero and ignoring writes to \texttt{rd\_addr == 0}.

\subsection{Read and write timing}

The register file has:

\begin{itemize}
\item two combinational read ports: \texttt{rs1\_data} and \texttt{rs2\_data} update when addresses change;
\item one synchronous write port: writes occur on the rising clock edge.
\end{itemize}
This is convenient for the single-cycle CPU because operands can be read, executed, and scheduled for write-back within one clock period. In the pipelined CPU, register-file timing interacts with hazards, forwarding, and write-back bypassing.

\bigskip
